% TOP_PROBLEM: {"id": 20001089, "title": "An explicit local-lemma certificate and bounded high-overlap neighbourhood for van der Waerden hypergraphs", "category_id": 2, "category": {"id": 2, "name": "combinatorics", "display_name": "Combinatorics", "description": "Counting problems, graph theory, discrete structures.", "slug": "combinatorics", "order_index": 2, "created_at": "2026-07-31T15:26:25.671Z"}, "status": "partially_solved", "problem_number": "AIM-COMBINATORICS-0214", "set_id": 14, "statusQualification": "Restores the exponent 2^{k^2} from the primary PDF typography."}
% FIRST_POSTED: 2026-10-01
% ENTRY_KIND: statement_only
% UnsolvedMath ID 20001089; source catalogue code AIM-COMBINATORICS-0214
% !TeX program = lualatex
% Definition and sources only; this file is not an LLM solution attempt.
\documentclass[11pt,a4paper]{article}
\usepackage[margin=25mm]{geometry}
\usepackage{fontspec}
\setmainfont{lmroman10-regular.otf}[BoldFont=lmroman10-bold.otf,
 ItalicFont=lmroman10-italic.otf,BoldItalicFont=lmroman10-bolditalic.otf]
\usepackage{amsmath,amssymb,mathtools}
\usepackage{unicode-math}
\setmathfont{latinmodern-math.otf}
\AtBeginDocument{\let\setminus\smallsetminus}
\usepackage{xurl,hyperref}
\hypersetup{colorlinks=true,urlcolor=blue}
\setcounter{secnumdepth}{0}
\setlength{\parindent}{0pt}
\setlength{\parskip}{5pt}
\setlength{\emergencystretch}{3em}
\begin{document}
\section*{An explicit local-lemma certificate and bounded high-overlap neighbourhood for van der Waerden hypergraphs}\label{20001089}
{\small UnsolvedMath ID 20001089 · AIM-COMBINATORICS-0214 · Combinatorics · AIM Workshop Problem Lists}
\subsection{Definitions and mathematical statement}
For an integer $k\ge2$, a nonconstant $k$-term arithmetic progression
in $[N]=\{1,\ldots,N\}$ is a set
\[
                   \{a,a+d,\ldots,a+(k-1)d\},
                 \qquad a,d\in\mathbb Z,\quad d\ge1,
\]
whose terms all lie in $[N]$. A two-coloring is any map
$c:[N]\to\{\mathrm{red},\mathrm{blue}\}$. A progression is
monochromatic when $c$ has the same value on all its terms.
Let $W(k)$ be the least $N$ for which every two-coloring of $[N]$
contains a monochromatic $k$-term progression. The van der Waerden
theorem ensures that $W(k)$ is finite.

The workshop asks for the true order of growth of $W(k)$ as
$k\to\infty$. It singles out the proposed upper bound
\[
                              W(k)\le2^{k^2}
                                      \qquad(k\ge2).
\]
The exponent is $k^2$; the expression is not $2k^2$, as a flattened
text extraction could suggest. The two-color convention is fixed
throughout and is not an additional growing parameter.
An upper bound must apply to every coloring of the entire interval,
while a lower bound is provided by a coloring of a shorter interval
avoiding all such monochromatic progressions. A coloring that avoids
only progressions with some specified differences does not suffice.
The broad growth question remains distinct from the particular
$2^{k^2}$ target: settling the latter would not itself determine
the sharp asymptotic order.
\subsection{Short English statement}
Determine the growth of the two-color van der Waerden numbers, including the proposed upper bound two to the power k squared.
\subsection{Sources}
\begin{itemize}
\item[S1] ULAM.AI and the UnsolvedMath Contributors, supplied problems.json, record 20001089 (AIM-COMBINATORICS-0214), AIM Workshop Problem Lists. Catalogue identity and source transcription; CC BY 4.0.
\url{https://huggingface.co/datasets/ulamai/UnsolvedMath}.
\item[S2] American Institute of Mathematics, Recent trends in additive combinatorics, item 1.8. Original workshop question underlying AIM-COMBINATORICS-0214.
\url{https://aimath.org/WWN/additivecomb/additivecomb.pdf}.
\end{itemize}
\end{document}
